\begin{table}[!htbp]

\caption{\label{tab:aggregation-bias-rent-bridge-2022}Aggregation bias with national quintiles and relative-expenditure weights, 2022}
\centering
\begin{tabular}[t]{lrrrrr}
\toprule
Country & Digit 2 gap & Digit 3 gap & Digit 4 gap & D2--D3 bias & D2--D4 bias\\
\midrule
DE & 1.72 & 0.16 & -- & 1.56 & --\\
FR & 0.25 & -0.38 & -- & 0.63 & --\\
ES & 0.84 & 0.39 & -- & 0.44 & --\\
NL & 0.42 & 2.62 & 3.38 & -2.20 & -2.96\\
BE & 0.44 & 1.36 & 1.94 & -0.93 & -1.51\\
\bottomrule
\end{tabular}

\begin{flushleft}
\footnotesize{\emph{Notes:} Quintiles are defined separately within each country using survey-weighted equivalised household income. For each product, its 2022 HICP weight is multiplied by the ratio of 2020 survey-weighted mean HBS expenditure in the quintile to the national mean; weights are then normalised within quintile over matched products. Available HBS actual and imputed rents (041 plus 042) are pooled and assigned the 041 HICP index and weight. The microdata contain no 042 variable for DE, FR or ES, so their rental cell uses actual rents only; the full bridge is available for NL and BE. This rental cell remains at digit 3 in the digit-4 calculation. Price reaggregation uses HICP weights. Annual item inflation is the mean of monthly year-on-year rates in 2022, requiring at least ten months. Gaps are Q1 minus Q5 in percentage points. COICOP digits 2--3 are available for DE, FR and ES, and digits 2--4 for NL and BE. Coverage can differ across digit depths.}
\end{flushleft}
\end{table}
